\documentclass[12pt, titlepage]{article}

\usepackage{amsmath, mathtools}

\usepackage[round]{natbib}
\usepackage{amsfonts}
\usepackage{amssymb}
\usepackage{graphicx}
\usepackage{colortbl}
\usepackage{xr}
\usepackage{hyperref}
\usepackage{longtable}
\usepackage{xfrac}
\usepackage{tabularx}
\usepackage{float}
\usepackage{siunitx}
\usepackage{booktabs}
\usepackage{multirow}
\usepackage[section]{placeins}
\usepackage{caption}
\usepackage{fullpage}

\hypersetup{
bookmarks=true,     % show bookmarks bar?
colorlinks=true,       % false: boxed links; true: colored links
linkcolor=red,          % color of internal links (change box color with linkbordercolor)
citecolor=blue,      % color of links to bibliography
filecolor=magenta,  % color of file links
urlcolor=cyan          % color of external links
}

\usepackage{array}

\externaldocument{../../SRS/SRS}



\input{../../Comments.text}
\input{../../Common.text}

\begin{document}

\title{Module Interface Specification for \progname{}}

\author{\authname}

\date{\today}

\maketitle

\pagenumbering{roman}

\section{Revision History}

\begin{tabularx}{\textwidth}{p{3cm}p{2cm}X}
\toprule {\bf Date} & {\bf Version} & {\bf Notes}\\
\midrule
Date 1 & 1.0 & Notes\\
Date 2 & 1.1 & Notes\\
\bottomrule
\end{tabularx}

~\newpage

\section{Symbols, Abbreviations and Acronyms}

See SRS Documentation \href{https://github.com/JackEaston4/Troitsky-Toolbox/blob/main/pdfs/SRS/SRS.pdf}{here}.

\wss{Also add any additional symbols, abbreviations or acronyms}

\newpage

\tableofcontents

\newpage

\pagenumbering{arabic}

\section{Introduction}

The following document details the Module Interface Specifications (MIS) for the Troitsky Toolbox,
a suite of services for the McMaster Troitsky Club, which fields teams to build bridges in the
national Troitsky competition. The prime service of this toolbox is the analysis of a bridge model.
The analysis draws upon materials testing results uploaded to the system in another service. Two other 
auxiliary services exist to facilitate club operations. The former is the knowledge repository to upload,
search for, preview and download documents. The latter is a login and role-based access control (RBAC) service,
which prevents unauthorized access of documents or use of services. This is necessary because the toolbox
will be acessible by a publicly-acessible website hosted on a web server.

Complementary documents include the System Requirement Specifications
and Module Guide.  The full documentation and implementation can be
found \href{https://github.com/JackEaston4/Troitsky-Toolbox}{here}.

\section{Notation}

\wss{You should describe your notation.  You can use what is below as
  a starting point.}

The structure of the MIS for modules comes from \citet{HoffmanAndStrooper1995},
with the addition that template modules have been adapted from
\cite{GhezziEtAl2003}.  The mathematical notation comes from Chapter 3 of
\citet{HoffmanAndStrooper1995}.  For instance, the symbol := is used for a
multiple assignment statement and conditional rules follow the form $(c_1
\Rightarrow r_1 | c_2 \Rightarrow r_2 | ... | c_n \Rightarrow r_n )$.

The following table summarizes the primitive data types used by \progname. 

\begin{center}
\renewcommand{\arraystretch}{1.2}
\noindent 
\begin{tabular}{l l p{7.5cm}} 
\toprule 
\textbf{Data Type} & \textbf{Notation} & \textbf{Description}\\ 
\midrule
character & char & a single symbol or digit\\
integer & $\mathbb{Z}$ & a number without a fractional component in (-$\infty$, $\infty$) \\
natural number & $\mathbb{N}$ & a number without a fractional component in [1, $\infty$) \\
real & $\mathbb{R}$ & any number in (-$\infty$, $\infty$)\\
string & $\mathbb{S}$ & an ordered array of characters (e.g. a string) \\
list & List<$E$> & an ordered collection of elements of type $E$ \\
\bottomrule
\end{tabular} 
\end{center}

\noindent
The specification of \progname \ uses some derived data types: sequences, strings, and
tuples. Sequences are lists filled with elements of the same data type. Strings
are sequences of characters. Tuples contain a list of values, potentially of
different types. In addition, \progname \ uses functions, which
are defined by the data types of their inputs and outputs. Local functions are
described by giving their type signature followed by their specification.

\section{Module Decomposition}

The following table is taken directly from the Module Guide document for this project.

\begin{table}[h!]
\centering
\begin{tabular}{p{0.3\textwidth} p{0.6\textwidth}}
\toprule
\textbf{Level 1} & \textbf{Level 2}\\
\midrule

{Hardware-Hiding} & ~ \\
\midrule

\multirow{7}{0.3\textwidth}{Behaviour-Hiding} & Input Parameters\\
& Output Format\\
& Output Verification\\
& Temperature ODEs\\
& Energy Equations\\ 
& Control Module\\
& Specification Parameters Module\\
\midrule

\multirow{3}{0.3\textwidth}{Software Decision} & {Sequence Data Structure}\\
& ODE Solver\\
& Plotting\\
\bottomrule

\end{tabular}
\caption{Module Hierarchy}
\label{TblMH}
\end{table}

\newpage
~\newpage

\section{MIS of Web Server} \label{WebServer} \wss{Use labels for
  cross-referencing}

\subsection{Module}

Web Server

\subsection{Uses}

Detect user input and display output to the user

\subsection{Syntax}

\subsubsection{Exported Constants}

None

\subsubsection{Exported Access Programs}

\begin{center}
\begin{tabular}{p{2cm} p{4cm} p{4cm} p{2cm}}
\hline
\textbf{Name} & \textbf{In} & \textbf{Out} & \textbf{Exceptions} \\
\hline
GetInputs & - & inputs: List<HtmlInput> & Server error, connection error  \\
\hline
DisplayOutput & HtmlElement, value: \textit{S} & Server error, connection error \\
\end{tabular}
\end{center}

\subsection{Semantics}

\subsubsection{State Variables}

\subsubsection{Environment Variables}

\begin{itemize}
  \item The HTML inputs (e.g. button clicks, text entered)
\end{itemize}

\subsubsection{Assumptions}

\begin{itemize}
  \item The user's browsers can display a HTML5 document and handle common inputs the document facilitates. This assumption is very reasonable because every modern browser has this capability.
\end{itemize}

\subsubsection{Access Routine Semantics}

Please note a rough sketch will be given in lieu of the actual implementation as implementation
varies by browser and device hardware. This is alright because this is an external component.

These functions do not exist directly. However, they summarize external activity.

\noindent GetInput()
\begin{itemize}
\item transition: poll each input-supporting element if any new input was given by the user 
\item output: a list of 2-tuples: (h: HtmlElement, value: $S$) 
\item exception: server error, connection error 
\end{itemize}

\noindent DisplayOutput(h: HtmlElement, value: $S$)
\begin{itemize}
\item transition: value is applied to the HTML element
\item output: none
\item exception: server error, connection error
\end{itemize}
\subsubsection{Local Functions}

None (varies by browser and device hardware)

\section{MIS of Access Control System} \label{ACS} 

\wss{You can reference SRS labels, such as R\ref{R_Inputs}.}

\wss{It is also possible to use \LaTeX for hypperlinks to external documents.}

\subsection{Module}

ACS

\subsection{Uses}

Log in, log out and permit/forbid an action a user wishes to perform.

\subsection{Syntax}

\subsubsection{Exported Constants}

None

\subsubsection{Exported Access Programs}

\begin{center}
\begin{tabular}{p{2cm} p{4cm} p{4cm} p{2cm}}
\hline
\textbf{Name} & \textbf{In} & \textbf{Out} & \textbf{Exceptions} \\
\hline
logIn & - & success: \textit{B}& - \\
\hline
logOut & - & success: -& \textit{B} \\
\hline
isAllowed & email, action & allowed: \textit{B}& - \\
\hline
setRole & email, role & success: \textit{B} & - \\
\hline
createAccount & email, role & success: \textit{B} & - \\
\hline
deleteAccount & email & success: \textit{B} & - \\
\hline
createRole & roleName & success: \textit{B} & - \\
\hline
setRoleName & oldRoleName, newRoleName & success: \textit{B} & - \\
\hline
setRoleActionPermit & roleName, action, allowed & success: \textit{B} & - \\
\hline
deleteRole & roleName & success: \textit{B} & - \\
\hline 

\end{tabular}
\end{center}

\subsection{Semantics}

\subsubsection{State Variables}

\begin{itemize}
  \item sessionAccount: the account of the current session
\end{itemize}

\subsubsection{Environment Variables}

\begin{itemize}
  \item connStr: \textit{S}. Login database connection string
\end{itemize}

\subsubsection{Assumptions}

None. If there is a database error it is logged. Furthermore, an error means the system is compromised.
Thus, as a fail-safe effort, the server
will automatically take itself down and alert the developers to prevent malicious access and control


\subsubsection{Access Routine Semantics}

\noindent logIn():
\begin{itemize}
  \item Find the email the session supplied and verified through SSO. If none found, return false.
  \item Find the account in the accounts table that matches the email and password
  \item If an account found, switch the session's account from none to the one found
  \item output: if one found, return true. Otherwise, return false.
\end{itemize}

\noindent logOut():
\begin{itemize}
  \item Switch the session's account to none
\end{itemize}

\noindent isAllowed(account, action):
\begin{itemize}
  \item Find the action id in the action table. If not found, return false.
  \item Find the account's role by querying the roles database. If not found, return false.
  \item Query the roles table to find the row with the row and action id. If not found, return false.
  \item Output: The value of the permitted column.
\end{itemize}

\noindent setRole(account, role):
\begin{itemize}
  \item Call isAllowed(sessionAccount, \textit{setRole}). If false, return false.
  \item Find the role in the roles table. If not found, return false.
  \item Invoke the accounts table to update the account's role.
  \item Output: return true.
\end{itemize}

\noindent createAccount(email role)
\begin{itemize}
  \item Call isAllowed(sessionAccount, \textit{createAccount}). If false, return false.
  \item If the role is not found in the roles table, return false.
  \item If the email is already in the accounts database, update the account's role.
  \item Else, add a new row to the accounts database, containing the email, and role.
  \item Return true.
\end{itemize}

\noindent deleteAccount(email)
\begin{itemize}
  \item Call isAllowed(sessionAccount, \textit{deleteAccount}). If false, return false.
  \item Delete the account from the accounts table.
  \item Return true
\end{itemize}

\noindent createRole(roleName)
\begin{itemize}
  \item Call isAllowed(sessionAccount, \textit{createRole}). If false, return false.
  \item If the role name is in the roles table, return false.
  \item Add the role to the roles table.
  \item Return true.
\end{itemize}

\noindent setRoleName(oldRoleName, newRoleName)
\begin{itemize}
  \item Call isAllowed(sessionAccount, \textit{setRoleName}). If false, return false.
  \item If the role name is in the roles table, return false.
  \item Set the role with the old name in the roles table to the new name.
  \item Return true
\end{itemize}

\noindent setRoleActionPermit(role, action, permitted)
\begin{itemize}
  \item Call isAllowed(sessionAccount, \textit{setRoleActionPermit}). If false, return false.
  \item If the role name is not in the roles table, return false.
  \item If the action name is not in the actions table, return false.
  \item Add or update a row in the action permit table, to the new values.
  \item Return true
\end{itemize}

\noindent deleteRole(role)
\begin{itemize}
  \item Call isAllowed(sessionAccount, \textit{deleteRole}). If false, return false.
  \item Check the accounts table. If any account has that role, return false.
  \item Delete the roles from the roles table. (Note: the delete will cascade to the action permit table)
  \item Return true.
\end{itemize}

\subsubsection{Local Functions}

\noindent init(): -> None
\begin{itemize}
  \item Connect to the database with the connection string. 
  \item If connection fails, stop server and alert developers.
\end{itemize}

\subsection{Database Tables}

Accounts: $(AccountID, Email, RoleID)$, primary key $AccountID$, foreign key $RoleID$, on delete restrict.
Role: $(RoleID, RoleName, RoleDescription)$, primary key $RoleID$
ActionPermit: $(RoleID, ActionID, Permitted)$, primary key $Permitted$, foreign keys $RoleID, ActionID$, both on delete cascade.
Actions: $(ActionID, ActionName, ActionDescription)$.

\section{MIS of Toolbox Broker} \label{Broker} 

\subsection{Module}

ToolboxBroker

\subsection{Uses}

Facilitate service-based architecture, keeping services separate and making the creation of new services easy.


\subsection{Syntax}

\subsubsection{Exported Constants}

None

\subsubsection{Exported Access Programs}

\begin{center}
\begin{tabular}{p{2cm} p{4cm} p{4cm} p{2cm}}
\hline
\textbf{Name} & \textbf{In} & \textbf{Out} & \textbf{Exceptions} \\
\hline
goToService & serviceName & success: $B$ & None \\
\hline
\end{tabular}
\end{center}

\subsection{Semantics}

\subsubsection{State Variables}

None

\subsubsection{Environment Variables}

None

\subsubsection{Assumptions}

None

\subsubsection{Access Routine Semantics}

\noindent goToService(serviceName):
\begin{itemize}
\item Find the service in the services table. If not found, tell web server to display service non-existant error and return false.
\item Find the list of actions in the services actions. 
\item For each action in the list, delete it if not permitted. Find permission status based on action permit table and current role.
\item If no actions left, display service unacessible through web server and return false.
\item Create new page for the service. On that page, list each action that is permitted. Then display it via web server.
\item Return true.
\end{itemize}

\subsubsection{Local Functions}

None necessary

\subsection{Database Tables}

Services: $(ServiceID, ServiceName, ServiceDescription)$, $ServiceID$ primary
ServiceActions: $(ServiceID, ActionID)$, both foreign to Services and Actions respectively, both on delete cascade.

\section{MIS of Materials Testing Manager} \label{MatlsTest} \wss{Use labels for
  cross-referencing}

\wss{You can reference SRS labels, such as R\ref{R_Inputs}.}

\wss{It is also possible to use \LaTeX for hypperlinks to external documents.}

\subsection{Module}

MatlsTest

\subsection{Uses}

Upload tests, let the uploader or an executive deem them unreliable


\subsection{Syntax}

\subsubsection{Exported Constants}

None

\subsubsection{Exported Access Programs}

\begin{center}
\begin{tabular}{p{2cm} p{4cm} p{4cm} p{2cm}}
\hline
\textbf{Name} & \textbf{In} & \textbf{Out} & \textbf{Exceptions} \\
\hline
submitTest & List<(materialName: \textit{S}, strength: \textit{R}) & - & - \\
\hline
markTest & testId: \textit{N}, isValid: \textit{B} & success: \textit{B} & - \\
\hline
parseTestSheet & - & List<(materialName, strength) & file invalid \\
\hline
getStrength & materialName: $S$ & strength: $I$ & no strength found \\
\end{tabular}
\end{center}

\subsection{Semantics}

\subsubsection{State Variables}

None. Memory stored in database (see below).

\subsubsection{Environment Variables}

\begin{itemize}
  \item testSheet - the CSV materials test sheet submitted 
\end{itemize}

\subsubsection{Assumptions}

None.

\subsubsection{Access Routine Semantics}

\noindent submitTest(List<(materialName: \textit{S}, strength: \textit{R})>)
\begin{itemize}
  \item For each element:
  \item \begin{itemize}
    \item Find the material id linked to the name in the materials table
    \item If material id found, add new row recording strength to materials strengths table
  \end{itemize}
\end{itemize}

\noindent markTest(testId: $I$, isValid: $B$)
\begin{itemize}
    \item Find the test with the id in the materials strength table.
    \item If not found, return false.
    \item Mark the text as $isValid$. Return true.
\end{itemize}

\noindent parseTestSheet()
\begin{itemize}
  \item Use a spreadsheet-style-document parser to parse the document into a List<(materialName, strength)>
  \item If parse failed, throw an exception.
  \item Otherwise, return the list.
\end{itemize}

\noindent getStrength(materialName)
\begin{itemize}
  \item Find the material id in the materials table. If not found, throw not found exception.
  \item Pick the most recent valid strength recorded (done through sorting by timestamp). Return it.
  \item If none found, throw not found exception.
\end{itemize}

\subsubsection{Local Functions}

None necessary.

\subsection{Database Tables}
\begin{itemize}
  \item Materials: $(MaterialID, MaterialName, MaterialDescription)$, $MaterialID$ primary.
  \item MaterialStrengths: $(MaterialID, Strength, Timestamp, AccountID, Valid)$, $(MaterialID, Timestamp)$ primary, both ids foreign. For material id, on delete cascade. For account id, on delete set null.
\end{itemize}

\newpage

\section{MIS of Knowledge Repositiory} \label{KnowRepo} \wss{Use labels for
  cross-referencing}

\subsection{Module}

KnowRepo

\subsection{Uses}

The knowledge repository. Upload, search, preview and download document.

\subsection{Syntax}

\subsubsection{Exported Constants}

None.

\subsubsection{Exported Access Programs}

\begin{center}
\begin{tabular}{p{2cm} p{4cm} p{4cm} p{2cm}}
\hline
\textbf{Name} & \textbf{In} & \textbf{Out} & \textbf{Exceptions} \\
\hline
uploadDoc & - & success: \textit{B} & - \\
\hline
searchDocs & searchQuery: $S$ & docIDs: List<$I$> & - \\
\hline
selectDoc & docID & success: $B$ - \\
\hline
alterDoc & docID & success: $B$ & - \\
\hline
deleteDoc & docID & success: $B$ & - \\
\end{tabular}
\end{center}

\subsection{Semantics}

\subsubsection{State Variables}

None.

\subsubsection{Environment Variables}

\begin{itemize}
  \item uploadedDoc - the document uploaded
  \item foundDoc - the document found
\end{itemize}

\subsubsection{Assumptions}

There is a shared drive which stores the documents. A database is inappropriate as the documents
are too vast.

\subsubsection{Access Routine Semantics}

\noindent uploadDoc()
\begin{itemize}
  \item Upload uploadedDoc to the drive.
  \item Output: Whether upload was successful
\end{itemize}

\noindent searchDoc(searchQuery)
\begin{itemize}
  \item Search the drive with the query
  \item Output: the list of document ids that match the query, sorted by relevancy, in decreasing order
\end{itemize}

\noindent selectDoc(docID)
\begin{itemize}
  \item Set foundDoc to the docID if the document exists.
  \item Returns: whether foundDoc was set.
  \item Note: The browser will handle preview and download
\end{itemize}

\noindent alterDoc(docID)
\begin{itemize}
  \item Replace the document at docID with the new the one at uploadedDoc.
  \item Returns: whether the replacement was successful
\end{itemize}

\noindent deleteDoc(docID)
\begin{itemize}
  \item Delete the document and docID
  \item Returns: whether the deletion was successful
\end{itemize}

\subsubsection{Local Functions}

\wss{As appropriate} \wss{These functions are for the purpose of specification.
  They are not necessarily something that is going to be implemented
  explicitly.  Even if they are implemented, they are not exported; they only
  have local scope.}

  None necessary.

\newpage

\section{MIS of Knowledge Transfer Manager} \label{KnowTrans} \wss{Use labels for
  cross-referencing}


\subsection{Module}

KnowTrans

\subsection{Uses}

Facilitate knowledge transfer between outgoing and incoming club officers.

\subsection{Syntax}

\subsubsection{Exported Constants}

None.

\subsubsection{Exported Access Programs}

\begin{center}
\begin{tabular}{p{2cm} p{4cm} p{4cm} p{2cm}}
\hline
\textbf{Name} & \textbf{In} & \textbf{Out} & \textbf{Exceptions} \\
\hline
prepareDocs & - & success: $B$ & - \\
\hline
retrieveDocs & position & List<docID>: $B$ & - \\
\hline
\end{tabular}
\end{center}

\subsection{Semantics}

\subsubsection{State Variables}

None

\subsubsection{Environment Variables}

\begin{itemize}
  \item uploadedDocs: the list of uploaded transition documents
\end{itemize}

\subsubsection{Assumptions}

None. \\

Please note the retrieveDocs will create a search and call the knowledge repository's search function.
Furthermore, an invalid position name will just yield no results.

\subsubsection{Access Routine Semantics}

\noindent prepareDocs():
\begin{itemize}
\item Figure out the user's position based on the role (check $AccountID$ of the session's user for $RoleID$)
\item Upload the documents to the drive under the correct tags
\item output: whether upload was successful
\item exception: \wss{if appropriate} 
\end{itemize}

\noindent retrieveDocs():
\begin{itemize}
  \item Figure out the user's position based on the role (check $AccountID$ of the session's user for $RoleID$)
  \item Prepare a query for transition docs. 
  \item Call KnowRepo's searchDocs() with the query.
  \item Returns: The result of serachDocs()
\end{itemize}
\subsubsection{Local Functions}

\noindent getPosition():
\begin{itemize}
  \item Query session's account, to get the $RoleID$
  \item Returns: The role ID
\end{itemize}



\newpage

\section{MIS of Model Analysis System} \label{ModAny} 

\subsection{Module}

ModAny

\subsection{Uses}

Do model analysis on an inputted step file.

\subsection{Syntax}

\subsubsection{Exported Constants}

None.

\subsubsection{Exported Access Programs}

\begin{center}
\begin{tabular}{p{2cm} p{4cm} p{4cm} p{2cm}}
\hline
\textbf{Name} & \textbf{In} & \textbf{Out} & \textbf{Exceptions} \\
\hline
analyzeModel & - & success: $B$ & - \\
\hline
\end{tabular}
\end{center}

\subsection{Semantics}

\subsubsection{State Variables}

\subsubsection{Environment Variables}

\begin{itemize}
  \item uploadedModel - the STEP model uploaded.
  \item modelVisualization - the visualization of the model displayed to the user
  \item modelAnalysisReport - the report generated from model analysis.
\end{itemize}

\subsubsection{Assumptions}

Failure is handled using the success boolean variable. Nonetheless, a large class getting extracted
model data will be developed when implementing this.

\subsubsection{Access Routine Semantics}

\noindent analyzeModel():
\begin{itemize}
\item Analyze model. Check its rules, whether the bridge will fail, and generate its score. 
\item Generate a document as a result, set the modelAnalysisReport variable to it. 
\item If there is a problem, return false.
\item Upload the document in modelAnalysisReport to the knowledge repository in the testing user's name
\item Returns: whether analysis was a success
\end{itemize}

\wss{A module without environment variables or state variables is unlikely to
  have a state transition.  In this case a state transition can only occur if
  the module is changing the state of another module.}

\wss{Modules rarely have both a transition and an output.  In most cases you
  will have one or the other.}

\subsubsection{Local Functions}

While the first three will be tested, these should not be otherwise used externally.

\noindent getBridgeFailure(bridgeData)
\begin{itemize}
  \item Run the bridge failure on each element (see VnV Plan, S. 5.2.1 for details)
  \item Returns: $true$ if bridge fails, else $false$.
\end{itemize}

\noindent isRulesCompliant(bridgeData)
\begin{itemize}
  \item Extract weight, length, width and height, and see if rules compliant (see VnV Plan, S. 5.2.2 for details)
  \item Returns: $true$ if bridge is rules compliant, else $false$.
\end{itemize}

\noindent calculateScore(hardConstraints: List<$I$>, scoreData<weight: $I$, score $I$>)
\begin{itemize}
  \item Calculate the score (see VnV Plan, S. 5.2.3 for details)
  \item Returns: the score ($I$).
\end{itemize}

The rest do not have simple unit tests, and probably will be private.

\noindent score(bridgeData)
\begin{itemize}
  \item Extract hard constraints, and list of (weight, score) tuples from bridge data.
  \item Call calculateScore with the output.
  \item Returns: the score that calculateScore returned
\end{itemize}

\noindent extractBridgeData()
\begin{itemize}
  \item Take the STEP file at uploadedModel. 
  \item With the aid of an external CAD library, obtain values for bridge data.
  \item Create a new BridgeData object and store the values in there.
  \item Returns: the BridgeData object, or null if there is an error
\end{itemize}

\noindent createVisualization()
\begin{itemize}
  \item Take the STEP file at uploadedModel. 
  \item With the aid of an external library, create the visualiztion of the model modelled.
  \item Store the model at modelAnalysisReport
  \item Returns: whether the visualization was a success
\end{itemize}

\noindent generateReport()
\begin{itemize}
  \item Create a summary document, containing details about the test.
  \item The document should contain at least: bridge failure, rules compliance, score, timestamp, testing user, and the bridge's visualization.
  \item Save the document to modelAnalysisReport
  \item Invoke KnowRepo's document, uploading the document to the knowledge repository
  \item Returns: whether the report generation and upload was a success 
\end{itemize}

\noindent displayReport()
\begin{itemize}
  \item Work with toolbox browser and web server to display the docuent just created.
\end{itemize}

\newpage

\newpage

\bibliographystyle {plainnat}
\bibliography {../../../refs/References}

\newpage

\newpage{}

\section{Appendix --- Generative AI Use Disclosure}


\subsection{AI Tools Used}

Gemini (version 3)

\subsection{How and Where Used}

Advisory only. 

\subsection{Verification of AI-Generated Content}

\wss{\begin{itemize}
    \item Describe how the team checked AI-generated or AI-assisted material
    before including it in the document.
    \item In particular, verify factual claims, technical assertions,
    requirements, calculations, references, and other information for which an
    AI system may produce incorrect or fabricated results.
    \item \textbf{The team is responsible for the accuracy, completeness,
consistency, and appropriateness of the submitted document, regardless of
whether material was generated by a human or an AI system.}
\end{itemize}}

All advice was critically thought of and all factual claims within the advice were 
verified before the advice was used.

\newpage{}

\section{Appendix --- Reflection}

\wss{Not required for CAS 741 projects}

The information in this section will be used to evaluate the team members on the
graduate attribute of Problem Analysis and Design.

\input{../../Reflection.text}

\begin{enumerate}
  \item What went well while writing this deliverable? 
  \item What pain points did you experience during this deliverable, and how
    did you resolve them?
  \item Which of your design decisions stemmed from speaking to your client(s)
  or a proxy (e.g. your peers, stakeholders, potential users)? For those that
  were not, why, and where did they come from?
  \item While creating the design doc, what parts of your other documents (e.g.
  requirements, hazard analysis, etc), it any, needed to be changed, and why?
  \item What are the limitations of your solution?  Put another way, given
  unlimited resources, what could you do to make the project better? (LO\_ProbSolutions)
  \item Give a brief overview of other design solutions you considered.  What
  are the benefits and tradeoffs of those other designs compared with the chosen
  design?  From all the potential options, why did you select the documented design?
  (LO\_Explores)
  \item (After you have implemented another team's module, which means this
  isn't filled in until after the original deadline). What did you learn by
  implementing another team's module? Were all the details you needed in the
  documentation, or did you need to make assumptions, or ask the other team
  questions?  If your team also had another team implement one of your modules,
  what was this experience like?  Are there things in your documentation you
  could have changed to make the process go more smoothly for when an
  ``outsider'' completes some of the implementation?

  Your reflection should be about concepts learned in your SE education, such as modular decomposition, cohesion, coupling, information hiding, etc.  You should critique the quality and utility of the specification you received.

  You should create a PR for your implemented module.  A link to the PR should be part of your reflection.

\end{enumerate}

\subsection{Michael}
\begin{itemize}
  \item I thought of how this would be implemented low-level and did it with time to spare.
  \item Motivation - sometimes I got distracted while writing this large document. Eventually I committed as it got later into the night.
  \item None. The earlier documents were well designed to facilitate the produciton of this one
  \item With unlimited resources, we could use premier external libraries instead of the free versions that will ultimately be used. Also, a premier firewall would be installed to prevent any attack.
  \item Another design I considered were stateful objects instead of delegating to databases. 
    While this may make a slight speedup during runtime, storage and durability could be a nightmare if
     JSON was used in lieu of a database. Also RBAC was considered being scrapped because most people are 
     trustworthy anyway, but security is worthwhile due to the occasional malicious outsider. 
     Lastly, I chose that materials tests are initially valid because the vast, vast majority will be.
\end{itemize}

\end{document}